\scanpage{29}
\begin{center}-- 214 --\end{center}
\hfill n$^\circ$ 391

son ordre est $h-1=17$. Elle est orthogonale à toutes les racines simples
sauf à $\alpha_1$ et l'on a $(\tilde\alpha|\alpha_1)=1$.

(V)\quad Ce qui a été vu au n$^\circ$ 10 montre que les \uline{coracines} de $R'$
s'obtiennent encore en remplaçant $\varepsilon_i$ par $\varepsilon_i^*$
dans les expressions des racines de $R'$, et $R'{}^*$ est donc isomorphe à $R'$.
Toutes les racines ayant même longueur, la \uline{forme bilinéaire canonique}
pour $R$ est $(x|y)/2h=(x|y)/36$ et la longueur des racines est
$1/3\sqrt2$, (\S\ 8, n$^\circ$ 5, \uline{Remarque}).
Le nombre des racines non orthogonales à une racine donnée est
$4h-6=66$ et le nombre $\gamma$ est égal à $h^2=2^2.3^4$.

(VI)\quad Le calcul des poids fondamentaux, formant la base duale de
$(\alpha_i^*)$, $(1\leq i\leq7)$, dans $V'$, donne :
\[
\omega_1=\varepsilon_8-\varepsilon_7
=2\alpha_1+2\alpha_2+3\alpha_3+4\alpha_4+3\alpha_5+2\alpha_6+\alpha_7
=\tilde\alpha
\]
\[
\begin{aligned}
\omega_2&=\frac12(\varepsilon_1+\varepsilon_2+\varepsilon_3+\varepsilon_4+
\varepsilon_5+\varepsilon_6-2\varepsilon_7+2\varepsilon_8)=\\
&=\frac12(4\alpha_1+7\alpha_2+8\alpha_3+12\alpha_4+9\alpha_5+6\alpha_6+3\alpha_7)
\end{aligned}
\]
\[
\begin{aligned}
\omega_3&=\frac12(-\varepsilon_1+\varepsilon_2+\varepsilon_3+\varepsilon_4+
\varepsilon_5+\varepsilon_6-3\varepsilon_7+3\varepsilon_8)=\\
&=3\alpha_1+4\alpha_2+6\alpha_3+8\alpha_4+6\alpha_5+4\alpha_6+2\alpha_7
\end{aligned}
\]
\[
\begin{aligned}
\omega_4&=\varepsilon_3+\varepsilon_4+\varepsilon_5+\varepsilon_6+
2(\varepsilon_8-\varepsilon_7)=\\
&=4\alpha_1+6\alpha_2+8\alpha_3+12\alpha_4+9\alpha_5+6\alpha_6+3\alpha_7
\end{aligned}
\]
\[
\begin{aligned}
\omega_5&=\frac12(2\varepsilon_4+2\varepsilon_5+2\varepsilon_6+
3(\varepsilon_8-\varepsilon_7))=\\
&=\frac12(6\alpha_1+9\alpha_2+12\alpha_3+18\alpha_4+15\alpha_5+10\alpha_6+5\alpha_7)
\end{aligned}
\]
\[
\omega_6=\varepsilon_5+\varepsilon_6-\varepsilon_7+\varepsilon_8
=2\alpha_1+3\alpha_2+4\alpha_3+6\alpha_4+5\alpha_5+4\alpha_6+2\alpha_7
\]
\[
\omega_7=\varepsilon_6+\frac12(\varepsilon_8-\varepsilon_7)
=\frac12(2\alpha_1+3\alpha_2+4\alpha_3+6\alpha_4+5\alpha_5+4\alpha_6+3\alpha_7).
\]
\mspage{28}
